\chapter{1947:Sir Robert Robinson}

\label{chap:chemistry1947}

The Nobel Prize in Chemistry for the year 1947 was awarded to Sir Robert Robinson.

\textbf{Motivation:}

for his investigations on plant products of biological importance, especially the alkaloids

\section{Sir Robert Robinson}

\subsection*{Early Life and Education}

Sir Robert Robinson was born on 1886-09-13 in Rufford, England.

\subsection*{Career and Major Research}

Robinson studied at the University of Manchester under William Henry Perkin Jr., and later held professorships at the universities of Sydney, Liverpool, and Manchester. In 1930 he became Waynflete Professor of Chemistry at Oxford, a post he held until 1955. He carried out landmark work on the structure and synthesis of natural products, including alkaloids such as morphine and strychnine, and of plant pigments such as the anthocyanins.

\subsection*{Nobel Prize-winning Work}

The work for which Sir Robert Robinson was awarded the Nobel Prize in Chemistry in 1947 was described by the Nobel Committee as: "for his investigations on plant products of biological importance, especially the alkaloids".

In 1917 Robinson achieved a one-step synthesis of tropinone, the key ring system of the alkaloids cocaine and atropine, from succinaldehyde, methylamine, and acetone. The synthesis was a striking demonstration of biomimetic thinking in organic chemistry.

\subsection*{Significance and Impact}

Robinson's work established the structures of major classes of plant products and showed how complex ring systems found in nature could be built in the laboratory. His electronic theory of organic reactions, presented in 1932, strongly influenced theoretical organic chemistry.

\subsection*{Later Career and Honors}

Robinson served as president of the Royal Society from 1945 to 1950. He was knighted in 1939 and was awarded the Order of Merit. He died in 1975.

\subsection*{Legacy}

Robinson's synthesis of tropinone remains a classic of organic synthesis, and his structural work on alkaloids and plant pigments is foundational to natural product chemistry.

\subsection*{Selected Publications}

\begin{itemize}

\item \emph{A Synthesis of Tropinone}, Journal of the Chemical Society, 1917.

\end{itemize}

\clearpage

\section*{Chapter Summary}

The 1947 prize recognized investigations of plant products of biological importance, especially the alkaloids. Robert Robinson's studies of alkaloid structures and his 1917 synthesis of tropinone laid the groundwork for the modern chemistry of natural products.
